\SectionSlide{Additional Details}

\begin{frame}{Evaluate Models Under the Intended Use}
  \begin{itemize}\setlength{\itemsep}{.4cm}
    \item Split by time or group when observations are dependent.
    \item Fit preprocessing and tune hyperparameters using training/validation data.
    \item Compare a simple baseline, a tree, a forest, and boosted trees under the same protocol.
    \item Select a metric for the decision: error, ranking, calibration, or cost.
    \item Measure operational impact separately from offline predictive accuracy.
  \end{itemize}
  \src{An OOB score is useful for suitable i.i.d. settings; it does not replace a time-aware or group-aware test.}
\end{frame}

\begin{frame}{A Newton Step as Weighted Least Squares}
  For $h_i>0$, complete the square:
  \[
    g_i f(x_i)+\frac12h_i f(x_i)^2
    =\frac12h_i\left(f(x_i)+\frac{g_i}{h_i}\right)^2
      -\frac{g_i^2}{2h_i}.
  \]
  \pause
  \vspace{.3cm}
  \begin{columns}[T,onlytextwidth]
    \column{.47\textwidth}
    \textbf{Pseudo-target}\par
    \[z_i=-\frac{g_i}{h_i}.\]
    \column{.47\textwidth}
    \textbf{Sample weight}\par
    \[h_i.\]
  \end{columns}
  \vspace{.25cm}
  Fit a tree to these weighted targets, with the tree regularizer added.\par
  \vspace{.2cm}
  For squared loss, $h_i=1$ and $z_i=y_i-s_i$: ordinary residuals.
\end{frame}

\begin{frame}{Multiclass Scores and Softmax}
  Let $s_c(x)$ be the score for class $c\in\{1,\ldots,C\}$.
  \[p_c(x)=\frac{e^{s_c(x)}}{\sum_{k=1}^{C}e^{s_k(x)}}.\]
  \[\ell(y,\boldsymbol s)=-\log p_y(x).\]
  \pause
  \vspace{.3cm}
  \begin{itemize}\spaceditems
    \item Trees provide real-valued score increments for the classes.
    \item A common formulation fits one scalar-output tree per class per round.
    \item Predict $\widehat y=\arg\max_c p_c(x)$.
  \end{itemize}
  \src{This is the usual class-wise scalar-tree view. Multiclass curvature and vector-leaf variants require additional details.}
\end{frame}

\begin{frame}{Efficient Split Search Uses Histograms}
  \begin{enumerate}\spaceditems
    \item Map each numerical feature into a set of bins.
    \item Accumulate $G$ and $H$ in each bin.
    \item Scan bin boundaries using prefix sums to evaluate split gain.
  \end{enumerate}
  \vspace{.2cm}
  \centering
  \begin{tikzpicture}
    \foreach \i/\hh in {0/.6,1/1.0,2/1.8,3/1.2,4/.8,5/1.6,6/.9,7/.4}{
      \draw[fill=bulletgreen!45,draw=forestgreen] (\i*.9,0) rectangle +(0.85,\hh);
    }
    \draw[slideorange,very thick] (3.55,-.2)--(3.55,2.0);
    \node[above,text=slideorange,font=\small] at (3.55,2.0) {candidate boundary};
  \end{tikzpicture}
  \vspace{.3cm}
  \par For missing values, compare adding their $G,H$ to the left or right child.
  \src{Schematic histogram.}
\end{frame}

\begin{frame}{OOB Coverage Across the Forest}
  For independent bootstrap samples across $B$ trees:
  \[q=\left(1-\frac1n\right)^m,\qquad
    K_i=|\mathcal B_i|\sim\operatorname{Binomial}(B,q).\]
  \[\E[K_i]=Bq,\qquad P(K_i=0)=(1-q)^B.\]
  \pause
  \vspace{.25cm}
  For large $n$, with $m=5n$ and $B=100$, each sample has about a $51\%$ probability of having no OOB tree.
  \vspace{.2cm}
  \begin{itemize}\spaceditems
    \item OOB error evaluates subensembles with varying numbers of trees.
    \item Repeated tuning against the OOB score can overfit the evaluation.
    \item Time dependence and grouped observations need suitable validation splits.
  \end{itemize}
\end{frame}

\begin{frame}{An Operational Case: Virtua Health, 2023}
  XGBoost ran hourly within the electronic medical record workflow at five hospitals.
  \vspace{.3cm}
  \centering
  \renewcommand{\arraystretch}{1.4}
  \begin{tabular}{lc}
    \toprule
    Reported operational metric & Value\\ \midrule
    XGBoost sensitivity & $91\%$\\
    XGBoost false-positive rate & $6\%$\\
    Existing commercial model's false-positive rate & $30\%$\\ \bottomrule
  \end{tabular}
  \vspace{.2cm}
  \par Three-week post-deployment evaluation: 1,896 patients, 98 sepsis cases.
  \vspace{.3cm}
  \par A matched-sensitivity comparison was not reported.\par
  Patient-outcome improvements were not established.
  \src{Optional discussion case.}
\end{frame}

